\documentclass[11pt]{article}
\usepackage[margin=0.85in]{geometry}
\usepackage[T1]{fontenc}
\usepackage{amsmath,amssymb,booktabs,longtable,microtype,xurl,hyperref}
\setlength{\emergencystretch}{3em}
\hypersetup{colorlinks=true,urlcolor=blue,linkcolor=blue}
\title{Failure Causes in the Current Collision-Avoidance Controller}
\author{Pan Dynamics Collision Avoidance Research}
\date{October 2, 2026}
\begin{document}
\maketitle

\section{Outcome criterion and scope}
The control objective is a collision-free encounter followed by recovery to
the given path and desired cruise speed. There is no prescribed recovery
deadline. Completing an eight-second observation window is not a success
criterion. The reporting utility now leads with collision-free recovery,
observed collisions, interruptions, and unconfirmed recovery; historical
window-completion fields remain only as experimental bookkeeping.

The production algorithm analyzed here is commit
\path{a29ccb6ff2856fe3df33efff4b1cb8619e523658}. Among its fourteen latest
extended runs, twelve complete collision-free nominal recovery. The
15 m/s turning-crossing run collides and subsequently recovers. The 8 m/s
braking-lead run is interrupted before recovery is confirmed. These last
two outcomes have different causes. In particular, the interrupted run is
not evidence of asymptotic CLF failure.

This investigation changes reporting only. The controller, CLF, target
model, constraints and execution rules in the working repository remain
unchanged. Instrumented replay and an equivalent geometric-solver
counterfactual run in isolated copies. No new executable backup or online
nonlinear acceptance layer is introduced.

\section{Turning crossing: the causal chain}
\subsection{Exact replay and initial seed}
All 31 controller calls from 0 to 1.50 s were replayed using the original
measured states, target epoch, held inputs and recursively updated controller
memory. Every returned first input matches the original record exactly.
Both PCBF and CLF stages run. Thus the analysis below concerns the actual
reported controller, rather than a merely similar cold-start problem.

The initial flow rollout has 3.16313 m minimum body separation on the
node/midpoint grid. The initial optimized controls, independently propagated
with ODE45, also remain separated: minimum 2.85540 m over the inspected
encounter. An immediately colliding cold flow seed is not the initiating
failure in this run.

\subsection{The local safe set is too optimistic}
At each hold, the controller shifts the preceding controls, rolls them
through the nonlinear model to construct an anchor, and solves one affine
PCBF problem followed by one CLF problem on that same anchor. A zero
PCBF optimum permits the CLF objective to use the entire local safe set.
The local safe set is built using frozen dynamics and ordinary-distance
multipliers; it is not a certified inner approximation of nonlinear safety.

\begin{center}\small
\begin{tabular}{rrrr}
\toprule
Hold time (s) & Affine minimum (m) & Nonlinear RK4 minimum (m) & PCBF optimum\\
\midrule
0.00 & 2.93720 & 2.85538 & $3.88\,10^{-9}$\\
0.75 & 0.11293 & 0.04565 & $2.33\,10^{-9}$\\
0.85 & 0.11232 & $-0.00855$ & $2.11\,10^{-8}$\\
0.90 & 0.12375 & 0.11023 & $1.06\,10^{-8}$\\
1.05 & 0.10000 & $-0.04018$ & $2.95\,10^{-9}$\\
1.10 & $-0.06815$ & $-0.07368$ & 0.10000\\
\bottomrule
\end{tabular}
\end{center}
The two minimum columns use matching node/midpoint times and the same
returned inputs. Positive values are Euclidean rectangle distance; negative
values are signed separating-axis gaps (penetration diagnostics), not the
controller's ordinary distance, which is zero during overlap. They are
predictions under fixed future inputs, not subsequent closed-loop execution.
The fresh flow at 0.90 s temporarily restores a separated prediction.
The permanent positive-slack sequence starts at 1.10 s.

The second stage is material. At 0.85 s, the raw PCBF solution has 0.09293 m
minimum nonlinear RK4 clearance; the CLF solution has $-0.00855$ m. At
1.05 s, those values are 0.00498 m and $-0.04018$ m. Both stages respect
the same local constraints. This does not justify skipping the CLF stage:
it identifies an inaccurate feasible set in which the CLF objective is
allowed to optimize.

At the decisive 1.05 s hold, compare the two predictions at 1.525 s:
\begin{itemize}
\item The minimum affine fixed-dual value is 0.099998997 m. Evaluating the
same multipliers on the exact rectangle at the affine pose changes this
value by less than $10^{-14}$ m; yaw convexification of the rectangle is
negligible at this particular point.
\item The nonlinear RK4 position differs from the affine prediction by
0.07057 m, and heading differs by 0.03139 rad, about $1.80^\circ$.
\item Holding the dual normal fixed, the position error removes 0.07043 m
of separation and the heading error removes another 0.07485 m. The exact
fixed-dual value on the nonlinear pose is $-0.0452861$ m. The independently
computed signed rectangle gap is $-0.04018$ m.
\end{itemize}
The 0.10 m buffer is therefore exhausted by errors in the predicted vehicle
response, even at a time explicitly represented by a collision constraint.

\subsection{Why the dynamics tangent overestimates steering response}
The CLF stage repeatedly uses the full first steering correction,
$\Delta\delta\simeq-0.075$ rad, about $-4.30^\circ$. A fixed trust box of
that size does not make the Fiala tire curve nearly linear throughout it.
At the initial hold, the front-force tangent around the flow anchor predicts
3225 N for the returned first input, while evaluating the actual Fiala curve
at that same state/input gives 1365 N. Releasing steering from a strongly
nonlinear tire operating point retains too much lateral force in the tangent
model. This pattern recurs while the CLF objective gradually reduces early
avoidance.

As another diagnostic, at the 1.05 s hold the state/input tangent to the
front tire force reaches 1.150 times the corresponding lateral capacity at
future time 1.15 s. These force numbers are an offline explanation of the
local Fiala approximation, not separately optimized force variables. The
actual nonlinear tire model itself remains bounded.

Writing the actual prediction as
\[
 e_{i+1}=A_i e_i+B_i\Delta u_i+R_i(e_i,\Delta u_i),
\]
the affine problem drops $R_i$. The fixed trust box and the 0.10 m node
buffer are not linked by a bound on the accumulated pose remainder. A zero
PCBF slack in that problem consequently need not imply a safe nonlinear
rollout. This is the initiating design defect in the measured collision.

\subsection{Intersample geometry and integration are different effects}
The 1.05 s affine prediction has 0.10000 m minimum distance at constrained
nodes/midpoints, but its densely sampled pose interpolant falls to
0.006911 m at 1.515 s. The time-aligned rotating rectangle geometry can
consume almost the entire node buffer between those samples. The same
returned controls under continuous ODE45 dynamics reach signed gap
$-0.086754$ m at 1.520 s. At 0.85 s the corresponding dense nonlinear
prediction reaches $-0.062909$ m. These are offline observations, not added
online checks.

This is not primarily RK4 integration error. At the common time 1.525 s,
the 1.05 s plan's one-step-per-hold nonlinear RK4 position differs from
ODE45 by only $8.15\,10^{-5}$ m, versus 0.07057 m between affine and RK4
predictions. Four and eight RK4 substeps reduce that integration difference
to $6.37\,10^{-7}$ and $5.03\,10^{-8}$ m. Differences between a sparse-grid
minimum and a dense minimum must not be misreported as integration error.
The intersample loss increases vulnerability, but the nonlinear prediction
already overlaps at a constrained midpoint.

\subsection{Overlap makes ordinary-distance restoration flat}
After the 1.05 s controls are shifted and rolled nonlinearly, the 1.10 s
anchor already overlaps the target in its future prefix. The ordinary
rectangle distance satisfies
\[
 d(\mathcal E,\mathcal T)=0
\]
throughout strict overlap. Its first-order variation provides no direction
that distinguishes shallow overlap from deeper overlap. The selected
anchor's dual normal norm is $9.67\,10^{-10}$. A fresh-flow probe at the
same measured state also overlaps at 1.525 s, with normal norm
$7.99\,10^{-10}$ and maximum row-Jacobian magnitude $1.96\,10^{-9}$.
Numerically the affected row becomes
\[
 0+0^T\Delta x\ge 0.10-\xi,
 \qquad\text{hence}\qquad \xi\ge0.10.
\]
The reported 0.10 slack is then an optimum of the relaxed local problem,
not a failed numerical search for that problem. Increasing or decreasing
penetration does not improve this row's objective. Every hold from 1.10
through 1.50 s tries fresh flow, but both anchors can share this degeneracy.
The finite positive-slack solution is still executable under the current
controller semantics.

The initialization need not be a safety certificate. However, a restoration
method that must repair overlapping initializations needs useful search
information there. Allowing arbitrary overlapping seeds while relying solely
on the ordinary distance gradient leaves this interface incomplete.

\subsection{Why the fresh flow does not repair the late anchor}
The target is propagated at the correct absolute times; target motion is
not omitted. The Gaussian guide uses nominal path progress, while the
bicycle follows a heading/speed field with its own transient dynamics and
seed-only input shaping. At the 1.10 s restart, the flow rollout enters the
circular exclusion envelope at 1.40 s. At 1.50 s its lateral position is
3.02530 m while the Gaussian requests 5.75143 m. Its course lags the field
direction by $19.64^\circ$, and its path station lags the nominal arrival
clock by 1.08140 m. The field feedback requests steering 0.37879 rad;
seed shaping produces 0.14386 rad. The later rectangle overlap is therefore
compatible with a correctly moving, geometrically sensible flow field.
The vehicle does not instantaneously follow that field.

This establishes a late-initialization mismatch, not a proof that no
physical avoidance trajectory exists from that measured state. Nor does it
show that the initial flow method was intrinsically inadequate: its initial
rollout was separated. The sequence is loss of reliable local prediction,
then an overlapping anchor, then flat restoration.

\subsection{Executed collision}
The actual closed-loop body contact is sampled from 1.511667 to 1.538333 s.
Maximum signed-axis penetration is 0.152024 m at 1.523333 s. At the 1.50 s
hold, the issued input is $[0.0693151,\;0.223441]^T$; both attempted PCBF
problems have optimum approximately 0.10 and the CLF solve also returns a
converged finite vector. This is a control-safety failure despite successful
local solves. The run later confirms nominal recovery at 12.00 s, so this
case is not a failure to return to nominal behavior.

\section{Braking lead: numerical interruption, with a controlled counterfactual}
The production run stops at 10.35 s. The offending geometry belongs to the
shifted anchor at relative prediction time 2.50 s, absolute time 12.85 s:
\begin{center}\small
\begin{tabular}{lr}
\toprule
Independent rectangle distance & 27.633427417605 m\\
Returned dual lower bound & 27.633425742001 m\\
Primal--dual difference & $1.675604\,10^{-6}$ m\\
Configured absolute comparison tolerance & $1.0\,10^{-6}$ m\\
\bottomrule
\end{tabular}
\end{center}
The distance subproblem reports exit flag $-7$, with the solver message
that the search direction is small before its requested residual tolerances
are met. Its multipliers are nonnegative and its normal obeys the unit-norm
bound. The returned bound is slightly conservative. Nevertheless, the
unconditional absolute-gap exception aborts controller formulation. This
happens before that hold's PCBF and CLF problems are solved. It does not
represent obstacle contact or infeasibility of the control problem.

The actual CLF values before interruption decrease through
$6.40941,\ 5.90169,\ 5.46242,\ 5.05962$ at holds 10.15--10.30 s, with CLF
slack of order $10^{-8}$. At interruption the lateral error is only
0.001089 m, but heading, speed and yaw-rate errors have not yet satisfied
the full recovery dwell.

For an isolated counterfactual, the same ordinary-distance dual optimum is
constructed from the exact closest-feature rectangle normal already
available in \texttt{predictiveSafetyGeometry.rectangle}. The nonnegative
multipliers are the positive/negative target-frame components of this
normal. During overlap the construction returns the same ordinary zero
distance and zero normal; it is not a signed-distance substitution. The
same primal--dual and norm conditions remain enforced. At the offending
pose its gap is within roundoff. No tolerance is loosened and no MPC or CLF
constraint/objective is changed.

Resuming the saved state and controller memory at 10.35 s with only this
geometric construction changed reaches the recovery set at 11.65 s and
sustains it to confirmation at 12.65 s. The minimum observed body distance
for the combined original and resumed trace is 0.138015427 m. Independent
Python geometry agrees and confirms the recovery dwell over 7843 samples.
The final transverse errors are
\[
 [-0.016428\,\mathrm m,\quad0.005097\,\mathrm{rad},\quad
 0.010695\,\mathrm{m/s},\quad-0.007850\,\mathrm{m/s},\quad
 -0.006839\,\mathrm{rad/s}]^T.
\]
This identifies the original interruption as a geometric numerical-interface
failure and demonstrates recovery with the same CLF in this counterfactual.
It is not an unqualified production pass: the production geometry has not
been changed in this task.

\section{Design implications and evidence limits}
The measured priorities are:
\begin{enumerate}
\item Make the local dynamics/safety approximation reliable for the
steering corrections actually used by the CLF stage. The approximation
error and available geometric margin need a mathematical relationship;
a fixed input box alone supplies none.
\item Give restoration meaningful search information when a seed overlaps.
The ordinary zero-distance plateau cannot measure penetration progress.
A flow seed is a useful starting point, but cannot by itself repair this
property of the optimization model.
\item Treat collision slack according to its actual meaning. An optimal
positive slack can describe an overlapping plan; numerical convergence is
not proof of nonlinear avoidance. The current relaxed execution behavior
is recorded explicitly, rather than silently redefined as hard safety.
\item Remove the dependency of control-problem availability on micrometer
optimality of a distant rectangle-distance solve. The demonstrated exact
geometric construction is a mathematical alternative for this subproblem,
not an additional controller.
\end{enumerate}
The evidence does not support removing CLF optimization, imposing a recovery
deadline, or attributing the braking interruption to CLF instability.
No production algorithm correction is claimed here.

Recovery is an empirical dwell test: lateral error 0.10 m, heading
$1^\circ$, speed error 0.10 m/s, lateral velocity 0.05 m/s and yaw rate
0.01 rad/s, held for 1 s. It is not a proof of asymptotic convergence. The
counterfactual has no minimum-time gate and stops upon this dwell; its
60 s observation cap was not reached and is not a convergence requirement.
A finite sampled simulation also does not certify infinite-future or
continuous-time physical safety. The compared isolated processes were not
used to make new controller runtime claims.

\section{Reproduction and artifact locations}
The compact machine-readable findings, geometry failure, exact-dual result,
source hashes and isolated-source diffs are in
\path{report/CONTROLLER_FAILURE_CAUSES_20261002/}. The raw instrumented
problems, controller snapshots, executed MATLAB/Python helpers, logs,
nonlinear prediction samples and continuation exports are retained in
\path{/home/zai/.cache/collisionAvoidance/failure-analysis-20261002/}
and copied to the external project archive at closure. The prior campaign
inputs are identified by hash in \texttt{provenance.json}.

The executed diagnostic helpers are \texttt{runFailureReplay},
\texttt{analyzeTurning}, \texttt{analyzeCritical}, \texttt{probeDynamics},
\texttt{integrationProbe}, and \texttt{continueBraking}, run using MATLAB
\texttt{-singleCompThread -batch}. ODE45 uses relative/absolute tolerances
$10^{-11}/10^{-12}$ and 31 output samples per hold. RK4 refinement keeps the
same controls and target epoch. All scenarios are deterministic; no random
seed or random perturbation is introduced. The initial diagnostic replay
wrapper hit a structure-collection error before its turning loop completed;
a cell-array collector fixed the wrapper, after which the full 31-call
exact replay completed. No controller change was needed for that wrapper
correction. This turn reruns targeted diagnostics and the report summarizer,
not the entire fourteen-case campaign or the prior unit-test suite.
\end{document}
